\documentclass[11pt]{article}
\usepackage{vinotheca}

\begin{document}

\vinothecaTitle{TasteRank Explorer}{Grape Reference}{One-Paragraph Profiles for 101 Varieties}
\vinothecaAuthor{Jure Skarabot}{May 2026}
\vinothecaCompanions{Summary \textperiodcentered{} Technical Appendix \textperiodcentered{} Methods Primer \textperiodcentered{} Data Appendix \textperiodcentered{} Grape Reference}

\section*{Introduction}
This reference provides a concise one-paragraph profile for each of the 101 grape varieties in the TasteRank database. Entries cover typical characteristics, principal growing regions, historical origins, and the variety's stylistic signature. Grapes are organised by colour (red, then white) and alphabetically within each section. For sensory profiles, community assignments, and TasteRank scores, see the companion Data Appendix.

\input{grape_entries.tex}

\section*{References}

\begin{itemize}[leftmargin=2em,itemsep=0.2em]
\item d'Agata, I. (2014). \emph{Native Wine Grapes of Italy}. University of California Press.
\item Galet, P. (2000). \emph{Dictionnaire encyclop\'edique des c\'epages}. Hachette, Paris.
\item Robinson, J., Harding, J. \& Vouillamoz, J. (2012). \emph{Wine Grapes: A Complete Guide to 1{,}368 Vine Varieties}. Ecco / HarperCollins.
\end{itemize}

\end{document}
